\cotefile{11}
\scanpage{29}
% L1–2 | top
Donc on trouve un foncteur \deleted{\unsure{contravariant}}
\[
\unsure{\mathcal V|V} \longrightarrow \text{Algèbres bivariantes jaugées}
\]

% L3–6 | upper
3) Dans $\mathcal V^{\unsure{o|0|\circ}}$ \deleted{\unsure{il y a un}} il y a objet initial : $k$.

$\unsure{\forall|\mathrm{tt}|\text{Si}}\, A,B\in\mathcal V^{o}$, il existe un coproduit dans $\mathcal V^{o}$ (\deleted{produit} ds $\mathcal V$) savoir le produit tensoriel gauche $A\otimes B$. Soit \deleted{\unsure{les}} $f^*:A\unsure{-|\to}A'$, $g^*:B-B'$, alors $f^*\otimes g^*:A\otimes B\to A'\otimes B'$

% L7 | upper middle
défini par
\[
f\otimes g^*(a\otimes b)=f^*(a)\otimes g^*(b)
\]

% L8–12 | upper middle, additions on both sides of the display
et il \unsure{leur|hom.} correspond\unsure{ent|ant} des coproduits. \deleted{\unsure{On les note}} On convient de mettre sur $A\otimes B$ l'augmentation définie par
\[
k_{A\otimes B}=k_A\otimes k_B
\]
(produit tensoriel \uline{gauche} d'homs homogènes d'algèbres graduées)
\margin{i.e. $k_{A\otimes B}(x\otimes y)=k_A(x)k_B(\unsure{x|y})(-1)^{Bx}$}
\deleted{\unsure{je dis que}} l'on \deleted{\unsure{ait}} alors \deleted{\unsure{On aurait}}
\[
\varphi(f^*\otimes g^*)=\varphi(f^*)\otimes\varphi(g^*)
\]
\margin{produit tensoriel \uline{gauche} d'homs. \uline{homogènes} des \deleted{\unsure{groupes|produits}} modules gradués}

% L13–16 | centre
i.e
\[
(f\times g)_*(a'\otimes b')=f_*(a')\otimes g_*(b')\,(-1)^{ga'} \tag*{$(\ast)$}
\]

\uline{Mais c'est faux} \deleted{si} car on ne suppose pas $f$ et $g$ de degrés \inserted{pairs}

Les deux mb\unsure{rs} de la formule $(\ast)$ diffèrent par un facteur
\[
(-1)^{ag+bf+Bf}
\]

% L17–19 | lower middle
Si dans la définition de $k_{A\otimes B}$ on n'avait pas mis de signe $(-1)^{Bx}$, le signe qui \unsure{rendrait|viendraient} \unsure{correct|comme} $\unsure{(\ast)}$ serait
\[
(-1)^{a'g+bf}
\]

% L20–22 | lower
Si \inserted{de plus} on se dispensait dans $\unsure{(\ast)|\ast}$ de prendre $f_*\otimes g_*$ au sens \deleted{\unsure{gradué|produits}} tordu, et qu'on \deleted{\unsure{prenne|prend}} enlève le \uline{signe} $(-1)^{ga'}$, on trouverait donc un signe
\[
(-1)^{bf}=(-1)^{B'f}\unsure{.}
\]

% L23–24 | lower
Si on prenait $k$ tordu, $\otimes$ ordinaire, le signe serait
\[
(-1)^{bf+(AB+A'B')}
\]

% L25–27 | foot, boxed formula and note at right
C'est donc le signe où on prend tout ordinaire ($k_{A\otimes B}\inserted{=k_A\otimes k_B}$ et $f_*\otimes g_*$) qui est le plus simple. Donc on trouve alors
\[
\boxed{(f\times g)_*=(f_*\overset{\mathrm{ord}}{\otimes}g_*)\bigl(\mathrm{id}_A\otimes[(-1)^B\unsure{\eta|D}_{B'}]^{\deg f}\bigr)}
\]
\margin{\deleted{\illegible{1}} Il n'y a aucun signe dès que \textcircled{$f$} est pair.}
% archive sheet number: 28

